\chapter{Форма перерізу}\label{ch:geo4}

\section{Параметри форми}\label{sec:geo4-params}

Форму межі плазми (LCFS) описують кількома числами, які читають за її крайніми точками. Позначимо $R_{\max}$, $R_{\min}$ --- крайні радіуси, $Z_{\max}$, $Z_{\min}$ --- крайні висоти, а $R_{Z_{\max}}$, $R_{Z_{\min}}$ --- радіуси найвищої й найнижчої точок.

\begin{outofcorpus}[означення~\cite{Manual}, розд.~4]
\begin{equation}\label{eq:geo4-shape}
\begin{gathered}
R_{\mathrm{geo}}=\frac{R_{\max}+R_{\min}}2,\qquad a=\frac{R_{\max}-R_{\min}}2,\qquad A=\frac{R_{\mathrm{geo}}}a,\qquad \epsilon=\frac1A,\\
\kappa=\frac{Z_{\max}-Z_{\min}}{2a},\qquad \dup=\frac{R_{\mathrm{geo}}-R_{Z_{\max}}}a,\qquad \dlo=\frac{R_{\mathrm{geo}}-R_{Z_{\min}}}a,\qquad \delta=\frac{\dup+\dlo}2 .
\end{gathered}
\end{equation}
$A$ --- аспектне відношення, $\epsilon$ --- обернене аспектне відношення, $\kappa$ --- витягнутість (elongation), $\delta$ --- трикутність (triangularity). $\delta>0$: верхівка зсунута \emph{всередину} від $R_{\mathrm{geo}}$ (класичний D); $\delta<0$ --- назовні («перевернутий D»).
\end{outofcorpus}

У коді це кілька рядків:

\code{geo/geom_tour.py}{37}{58}{(параметри форми, площа, об'єм)}

\noindent Рядки 40--46 --- дослівно~\eqref{eq:geo4-shape}. \texttt{argmax}/\texttt{argmin} по $Z$ знаходять верхню й нижню точки. Похибка --- порядку відстані між точками контуру: для EFIT-контуру це частки сантиметра, для щільно заданої аналітичної кривої вона нехтовно мала. Рядки 53--58 --- площа перерізу за формулою Гаусса («шнурка») і об'єм тора за теоремою Паппа--Гульдіна $V=2\pi R_{\mathrm c}S$, де $R_{\mathrm c}$ --- радіус центроїда перерізу, \emph{не} $R_{\mathrm{geo}}$.

\begin{figure}[!tb]\centering
\begin{minipage}[c]{0.44\textwidth}\centering
\includegraphics[width=\textwidth]{geo_diiid_shape.jpg}
\end{minipage}\hfill
\begin{minipage}[c]{0.53\textwidth}
\caption{Геометрія DIII-D \#203702, кадр 75: поверхні $\psiN=0{,}1\ldots0{,}9$ (сірі), LCFS з набору EFIT (синя), відкриті поверхні $\psiN>1$ (штрих), магнітна вісь, X-точка і точки, за якими читають параметри~\eqref{eq:geo4-shape}: крайні радіуси (квадрати), найвища й найнижча точки (трикутники). Найнижча точка межі збігається з X-точкою. Вісь зсунута від $R_{\mathrm{geo}}$ назовні (зсув Шафранова) і вгору від $Z_{\mathrm{geo}}$, бо X-точка «відтягує» низ межі. Власний розрахунок, код: \texttt{geometry/code/geom\_tour.py}; дані --- Fusion Equilibrium Challenge, CC BY 4.0.}\label{fig:geo4-diiid}
\end{minipage}
\end{figure}

\begin{table}[!tb]\centering\small
\caption{Параметри форми двох прикладів (запуск \texttt{geom\_tour.py})}\label{tab:geo4-shape}
\begin{tabular}{@{}lcc@{}}\toprule
& DIII-D \#203702, кадр 75 & JET 1975 (межа Міллера)\\\midrule
$R_{\mathrm{geo}}$, м & 1,656 & 2,960\\
$a$, м & 0,605 & 1,250\\
$A=R_{\mathrm{geo}}/a$ & 2,74 & 2,37\\
$\kappa$ & 1,70 & 1,68\\
$\dup$ / $\dlo$ & 0,38 / 0,69 & 0,348 / 0,348\\
$Z_{\mathrm{geo}}$, м & $-0{,}070$ & 0\\
площа перерізу $S$, м$^2$ & 1,77 & 8,12\\
радіус центроїда $R_{\mathrm c}$, м & 1,591 & 2,852\\
об'єм $V$, м$^3$ & 17,7 & 145,5\\
конфігурація & діверторна, нижня X-точка & лімітерна\\\bottomrule
\end{tabular}
\end{table}

Порівняння в табл.~\ref{tab:geo4-shape} повчальне: за формою перерізу DIII-D і JET дуже схожі ($A\approx2{,}4$--$2{,}7$, $\kappa\approx1{,}7$), а за розміром JET удвічі більший лінійно й у $\approx8$ разів за об'ємом. Відмінність форми --- у трикутності: DIII-D асиметричний ($\dlo\approx2\dup$), бо низ межі проходить через X-точку. JET 1975 симетричний, бо його межу задано аналітично (п.~\ref{sec:geo4-jet}).

\section{Форма Міллера}\label{sec:geo4-miller}

Щоб \emph{задати} межу кількома параметрами (а не прочитати з EFIT), найчастіше беруть параметризацію Міллера~\cite{Miller1998}:
\begin{equation}\label{eq:geo4-miller}
R(t)=R_0+a\cos\bigl(t+\arcsin\delta\,\sin t\bigr),\qquad Z(t)=\kappa a\sin t,\qquad t\in[0,2\pi).
\end{equation}
Перевірка: на $t=0$ і $\pi$ маємо $R=R_0\pm a$, тобто $R_{\mathrm{geo}}=R_0$ і півширина $a$. На $t=\pi/2$ маємо $Z=\kappa a$ і $R=R_0+a\cos(\pi/2+\arcsin\delta)=R_0-a\delta$. Отже, $\kappa$ і $\delta$ у~\eqref{eq:geo4-miller} --- \emph{точно} ті самі, що в~\eqref{eq:geo4-shape}; для JET скрипт відтворює їх до останнього знака. Повна модель Міллера містить ще радіальні похідні (зсув Шафранова $\pa R_0/\pa r$, шир витягнутості й трикутності), але для однієї межі досить~\eqref{eq:geo4-miller}.

\begin{figure}[!tb]\centering
\includegraphics[width=\textwidth]{geo_miller.jpg}
\caption{Форма Міллера~\eqref{eq:geo4-miller} при $R_0=\SI{3}{м}$, $a=\SI{1}{м}$. Ліворуч: витягнутість при $\delta=0$ (еліпси). У центрі й праворуч: додатна й від'ємна трикутність при $\kappa=1{,}7$. Власний розрахунок, функція \texttt{tokviz.surfaces.d\_shape}.}\label{fig:geo4-miller}
\end{figure}

\code{viz/src/tokviz/surfaces.py}{150}{166}{(D-переріз Міллера)}

\noindent Зверніть увагу на застереження в рядках 158--160: D Міллера --- \emph{не} та «D принстонського типу», якої набуває котушка TF. Котушку визначає умова чистого натягу (п.~\ref{sec:geo6-coil}), а не параметри форми. Схожість букв --- збіг: плазмі D-форму дають котушки полоїдального поля, котушці --- магнітні сили.

\section{Навіщо плазмі форма}\label{sec:geo4-why}

\paragraph{Більший струм за того самого $q$.} Для еліптичного перерізу у циліндричному наближенні (стандартна оцінка, поза корпусом)
\begin{equation}\label{eq:geo4-qcyl}
q_{\mathrm{cyl}}\approx\frac{2\pi a^2\Bzero}{\mu_0R_0\Ip}\cdot\frac{1+\kappa^2}2 .
\end{equation}
При $\kappa=1{,}7$ множник $(1+\kappa^2)/2=1{,}95$: за тієї самої межі стійкості $q$ плазма може нести майже вдвічі більший струм, а утримання в токамаках зростає зі струмом. Проєкт JET оцінював цей виграш у $1{,}7$ раза~\cite[розд.~5]{Manual}. Точніша формула для D-перерізу з тороїдальною поправкою --- \cite[розд.~5, (5.4)]{Manual}; вона дає для JET $q(a)=5{,}99$ (розд.~\ref{ch:geo6}).

\paragraph{Ціна: вертикальна нестійкість.} Щоб витягнути плазму, потрібне зовнішнє поле з від'ємним індексом спаду; тоді плазма вертикально нестійка і потребує пасивної стінки та активного керування~\cite[розд.~4]{Manual}.

\paragraph{Трикутність і дивертор.} За стандартною теорією (поза корпусом) трикутність підвищує межі стійкості за тиском, а X-точка на кінці «трикутника» дає дивертор. Тому в DIII-D $\dlo$ більша, ніж $\dup$. Від'ємна трикутність --- окремий напрям досліджень~\cite[розд.~4]{Manual}.

\paragraph{Заповнення камери.} Кругла плазма в D-подібній камері лишає порожніми «кути». За тієї самої площі поверхні D-переріз вміщує більший об'єм; саме так аргументували форму JET~\cite[розд.~5]{Manual}.

\section{Як задали межу JET 1975}\label{sec:geo4-jet}

Проєкт JET дає $R_0=\SI{2,96}{м}$, $a=\SI{1,25}{м}$ і напіввисоту $b=\SI{2,10}{м}$ (тобто $\kappa=b/a=1{,}68$), але не дає трикутності. Натомість у ньому є таблиця 30 точок межі у верхній зовнішній чверті~\cite[с.~332]{JET1975}. Проєкт~\cite{ProjJET} підганяє до неї форму Міллера і бере з підгонки лише $\delta$:

\code{viz/src/tokviz/machine/components/plasma.py}{69}{82}{(підгонка трикутності до таблиці межі)}

\noindent Підгонка $(a_t,\kappa_t,\delta)$ методом Нелдера--Міда за середньоквадратичною відстанню від 30 точок таблиці до кривої Міллера (рядки 72--77). Крива на внутрішньому екваторі проходить через край лімітера $R=\SI{1,71}{м}$: центр \texttt{r\_in + a} (рядок 76). Рядки 79--80: два проходи з меншим кроком симплекса.

\code{viz/src/tokviz/machine/components/plasma.py}{91}{118}{(межа плазми JET)}

\noindent Рядки 99--103: з підгонки береться лише $\delta$ (округлена до $10^{-4}$), а $R_0$, $a$, $\kappa$ --- з таблиці основних параметрів проєкту. Рядки 104--113: перевірка, що плазма не перетинає стінку камери. Якщо зазор менший за допустимий, $|\delta|$ зменшують кроками. Рядки 114--115: зазори на екваторі. Результат (маніфест \texttt{data/bake\_jet1975/manifest.json}): $\delta=0{,}3476$ при середньоквадратичному відхиленні від таблиці \SI{1,7}{см}; сама підгонка дає $\kappa_t=1{,}745$, тобто таблиця витягнутіша за $b/a$. Зазори до стінки на екваторі --- \SI{17,9}{см} ззовні й \SI{4,6}{см} всередині, найменший зазор по всьому периметру --- \SI{3,5}{см}, обмеження $\delta$ не спрацювало. Сама таблиця суперечлива: її верхня точка на 9,5~см вища за $b$ і виходить за стінку, тому таблицю використано лише як \emph{форму}, а не як межу.

\begin{practicum}[форма в часі]
(1)~Для всіх кадрів DIII-D \#203702 з \texttt{efit\_lcfs\_r/z} порахуйте \texttt{shape\_params} і побудуйте $\kappa(t)$, $\dup(t)$, $\dlo(t)$. Коли з'являється X-точка (стрибок $\dlo$)? (2)~Те саме для MAST \texttt{mast\_shot\_28348}: яке аспектне відношення має сферичний токамак? (3)~Для межі Міллера JET порівняйте $S$ і $V$ з еліпсом тих самих $a$, $\kappa$ (див. задачу~\ref{pr:geo7-pappus}).
\end{practicum}
